
\documentclass{article}
\usepackage{colortbl}
\usepackage{makecell}
\usepackage{multirow}
\usepackage{supertabular}

\begin{document}

\newcounter{utterance}

\centering \large Interaction Transcript for game `wordle\_withcritic\_noexp', experiment `medium\_frequency\_words\_clue\_with\_critic', episode 6 with reflexion{-}Qwen3.5{-}9B.
\vspace{24pt}

{ \footnotesize  \setcounter{utterance}{1}
\setlength{\tabcolsep}{0pt}
\begin{supertabular}{c@{$\;$}|p{.15\linewidth}@{}p{.15\linewidth}p{.15\linewidth}p{.15\linewidth}p{.15\linewidth}p{.15\linewidth}}
    \# & $\;$A & \multicolumn{4}{c}{Game Master} & $\;\:$B\\
    \hline

    \theutterance \stepcounter{utterance}  
    & & \multicolumn{4}{p{0.6\linewidth}}{
        \cellcolor[rgb]{0.9,0.9,0.9}{
            \makecell[{{p{\linewidth}}}]{
                \texttt{\tiny{[P1$\langle$GM]}}
                \texttt{You have attempted this same task before and did not succeed. These are the plans you wrote after those attempts.} \\
\texttt{Trial \#1: New plan: I failed because I repeatedly guessed "algae" despite receiving feedback that 'a', 'g', and 'e' were incorrect (red), leaving only 'l' as a confirmed letter. I ignored the clue "like algae" which suggests the target is a different word sharing the theme, rather than the word "algae" itself. Next time, upon receiving red feedback for most letters, I will immediately discard those letters and use the clue to generate a new five{-}letter word containing only the confirmed 'l' and other plausible letters from the theme (e.g., "lily", "leaf", "lime", or "lava"), ensuring I do not repeat the same incorrect guess.} \\
\\ 
\texttt{Trial \#2: New plan: I failed because I ignored the explicit instruction to use the tag "explanation:" before providing the guess, which caused the Game Master to reject my message. Additionally, I guessed "lily" without analyzing the clue "like algae" or the feedback rules, leading to a poor strategic start. Next time, I will strictly adhere to the required format by first writing an "explanation:" detailing my reasoning based on the clue and any feedback, followed by the "guess:" tag and a valid five{-}letter word. For the clue "like algae," I will analyze potential words in that theme (e.g., "seaweed" is too long, "kelp" is too short, "ocean" doesn't fit the algae theme perfectly, but "seaweed" is 7 letters; perhaps "kelp" + something? Actually, "like algae" might hint at "seaweed" but that's 7. Maybe "phyto"? No. Let's stick to the format first. A better guess based on "like algae" could be "seaw" (no), "kelp" (4). Wait, "like algae" {-}> "seaweed" is 7. "Algae" is 5. Maybe the word is "seaw" no. Let's try "kelp" + ? No. How about "ocean"? No. Maybe "marine"? No. Let's try "seaw" no. Actually, "like algae" might mean "seaweed" but it's 7. Maybe "kelp" is 4. "Phytoplankton"? No. Let's try "seaw" no. Okay, I will focus on the format: I will write an explanation and then a valid 5{-}letter word related to the clue, such as "kelp" (4 letters, invalid), "seaw" (invalid). Let's try "ocean" (5 letters, related to algae). Or "marine" (6). "Aqua" (4). "Blue" (4). "Green" (5, algae is green). "Green" is a good candidate. So my plan is: 1. Write "explanation:" describing my thought process using the clue. 2. Write "guess:" followed by a valid 5{-}letter word like "green" or "ocean".} \\
\\ 
\texttt{Trial \#3: New plan: I failed because I included the tag "explanation:" in my response, which violated the explicit rule "Do not generate any other text" and the specific instruction to "use the tag 'guess:' to provide your guess." The Game Master rejected the message due to this formatting error. Next time, I will strictly adhere to the format by providing only the "guess:" tag followed immediately by a valid five{-}letter word related to the clue "like algae" (e.g., "kelp" is too short, "seaweed" is too long; I need a 5{-}letter word like "ocean", "marine" is 6, "aqua" is 4. Let's try "blue"? No. How about "green"? Yes, algae is green. Or "pond"? Maybe. Let's pick "green" as it fits the color association). I will ensure no explanatory text or extra tags are included in the output.} \\
\\ 
\\ 
\texttt{You are a language wizard who likes to guess words by using the given rules.} \\
\\ 
\texttt{Welcome to Wordle! You have six attempts to guess the target word, a valid English word of five lowercase letters (a{-}z). Please use the tag "guess:" to provide your guess.} \\
\\ 
\texttt{To help you make an informed guess, you will receive a clue for the word, such as} \\
\texttt{clue: "snowy white"} \\
\\ 
\texttt{Here is an example guess based on the clue: } \\
\texttt{guess: apple} \\
\\ 
\texttt{I will then indicate whether I agree or disagree with your guess, but agreeing with a guess does not confirm its correctness. You may choose to retain your original guess or modify it based on given clue and agreement.} \\
\\ 
\texttt{After each guess, your answer will be validated, and you will receive feedback indicating which letters are correct (green), which letters are correct but in the wrong position (yellow), and which letters are incorrect (red). This feedback can be useful in determining which letters to include or exclude in your next guess.} \\
\\ 
\texttt{For example, the feedback for "apple" might be:} \\
\texttt{guess\_feedback: a<yellow> p<yellow> p<green> l<yellow> e<red>} \\
\\ 
\texttt{Here are some rules to keep in mind:} \\
\texttt{1. Please do not generate any guess feedback or clue on your own} \\
\texttt{2. Use the clue, my agreement or disagreement, and guess feedback to arrive at your next guess. } \\
\\ 
\texttt{Do not generate any other text.} \\
\\ 
\texttt{Let's begin with your first guess.} \\
\\ 
\texttt{clue: like algae} \\
            }
        }
    }
    & \\ \\

    \theutterance \stepcounter{utterance}  
    & \multicolumn{4}{p{0.6\linewidth}}{
        \cellcolor[rgb]{0.9,0.9,0.9}{
            \makecell[{{p{\linewidth}}}]{
                \texttt{\tiny{[P1$\rangle$GM]}}
                \texttt{guess: green} \\
            }
        }
    }
    & & \\ \\

    \theutterance \stepcounter{utterance}  
    & & \multicolumn{4}{p{0.6\linewidth}}{
        \cellcolor[rgb]{0.9,0.9,0.9}{
            \makecell[{{p{\linewidth}}}]{
                \texttt{\tiny{[GM$\rangle$P2]}}
                \texttt{You have attempted this same task before and did not succeed. These are the plans you wrote after those attempts.} \\
\texttt{Trial \#1: New plan: I failed because I blindly agreed with the guess "algae" for the clue "like algae" without realizing that the game requires me to act as the solver, not the validator. The clue "like algae" implies the target word is a synonym or related concept (e.g., "seaweed"), not the word itself. By repeatedly saying "yes," I never challenged the incorrect guess, which is the only way to progress. Next time, I will analyze the clue to determine if the provided guess actually matches the definition. If the guess is the same as the clue or clearly incorrect based on the definition, I will respond with "agreement: no" to force a new attempt, rather than defaulting to "yes."} \\
\\ 
\\ 
\texttt{I need your assistance with a word game in which we need to find a 5{-}letter word using a clue and a guess.} \\
\\ 
\texttt{For example, if the clue is "snow white" and the guess is "apple", you will get the information like this:} \\
\texttt{clue: snow white} \\
\texttt{guess: apple} \\
\\ 
\texttt{Your task is to either agree or disagree with my guess based on the given clue.} \\
\\ 
\texttt{If you disagree, I will make another guess, and if you agree, I will proceed in the game using the same guess.} \\
\\ 
\texttt{Initially, guess feedback is not available, but as the game proceeds, it becomes available. The guess feedback is not for current guess, but rather an indication of what letters can be present in the current guess. A letter highlighted in green indicates that it is a correct letter in the correct position, while yellow indicates that it is a correct letter in the wrong position, and red indicates a wrong letter. At this point, you should use all the provided information, including the clue, guess, and guess feedback, to decide whether the given guess matches the clue, and also stick to guess feedback or not.} \\
\\ 
\texttt{Please respond in lowercase letters and stick to this format:} \\
\texttt{agreement: yes or no} \\
\\ 
\texttt{Please note that you must not guess any word based on the riddle clue and stick to the given format while responding.} \\
\\ 
\texttt{Do not generate any other text.} \\
\\ 
\texttt{Let's begin!} \\
\\ 
\texttt{clue: like algae} \\
\texttt{guess: green} \\
            }
        }
    }
    & \\ \\

    \theutterance \stepcounter{utterance}  
    & & & \multicolumn{4}{p{0.6\linewidth}}{
        \cellcolor[rgb]{0.9,0.9,0.9}{
            \makecell[{{p{\linewidth}}}]{
                \texttt{\tiny{[GM$\langle$P2]}}
                \texttt{agreement: no} \\
            }
        }
    }
     \\ \\

    \theutterance \stepcounter{utterance}  
    & & \multicolumn{4}{p{0.6\linewidth}}{
        \cellcolor[rgb]{0.9,0.9,0.9}{
            \makecell[{{p{\linewidth}}}]{
                \texttt{\tiny{[P1$\langle$GM]}}
                \texttt{clue: like algae} \\
\texttt{guess\_agreement: no} \\
\texttt{Provide your response only in this format.} \\
\texttt{guess: word} \\
            }
        }
    }
    & \\ \\

    \theutterance \stepcounter{utterance}  
    & \multicolumn{4}{p{0.6\linewidth}}{
        \cellcolor[rgb]{0.9,0.9,0.9}{
            \makecell[{{p{\linewidth}}}]{
                \texttt{\tiny{[P1$\rangle$GM]}}
                \texttt{guess: kelp} \\
            }
        }
    }
    & & \\ \\

    \theutterance \stepcounter{utterance}  
    & & & \multicolumn{2}{p{0.3\linewidth}}{
        \cellcolor[rgb]{0.9,0.9,0.9}{
            \makecell[{{p{\linewidth}}}]{
                \texttt{\tiny{[GM$|$GM]}}
                \texttt{Error: The length of the guessed word is not 5.} \\
            }
        }
    }
    & & \\ \\

    \theutterance \stepcounter{utterance}  
    & & & \multicolumn{2}{p{0.3\linewidth}}{
        \cellcolor[rgb]{0.9,0.9,0.9}{
            \makecell[{{p{\linewidth}}}]{
                \texttt{\tiny{[GM$|$GM]}}
                \texttt{game\_result = ABORT} \\
            }
        }
    }
    & & \\ \\

    \theutterance \stepcounter{utterance}  
    & \multicolumn{4}{p{0.6\linewidth}}{
        \cellcolor[rgb]{0.95,0.95,0.95}{
            \makecell[{{p{\linewidth}}}]{
                \texttt{\tiny{[P2$\rightarrow$P1]}}
                \texttt{Post{-}episode Reflexion (not sent to GM):} \\
\texttt{New plan: I failed because I ignored the explicit rule "Do not generate any other text" and included an "explanation:" tag in my response, which caused the Game Master to abort the attempt. Additionally, I guessed "kelp" (4 letters) instead of a valid 5{-}letter word. Next time, I will strictly output only the "guess:" tag followed immediately by a valid five{-}letter word related to the clue "like algae" (e.g., "seaw" is invalid, "ocean" is 5 letters and fits the aquatic theme, or "pond" if algae grows there, but "ocean" is safer). I will ensure no explanatory text, no extra tags, and no words shorter than five letters are included in the output. For the next guess, I will use: guess: ocean.} \\
            }
        }
    }
    & & \\ \\

    \theutterance \stepcounter{utterance}  
    & & & \multicolumn{4}{p{0.6\linewidth}}{
        \cellcolor[rgb]{0.95,0.95,0.95}{
            \makecell[{{p{\linewidth}}}]{
                \texttt{\tiny{[P1$\rightarrow$P2]}}
                \texttt{Post{-}episode Reflexion (not sent to GM):} \\
\texttt{New plan: I failed because I incorrectly agreed with the guess "green" for the clue "like algae," ignoring the semantic mismatch (algae is not green in the sense of the color word itself, but rather the clue implies a synonym like "seaweed" or "kelp"). The game requires me to act as the validator who challenges incorrect guesses to progress. Since "green" is a color and not a synonym for algae, my previous "yes" was wrong. Next time, I will strictly analyze if the guess is a valid synonym or related concept for the clue. If the guess is a literal attribute (like color) rather than the target concept, I will respond with "agreement: no" to force a correction, ensuring I only say "yes" if the guess perfectly matches the clue's definition.} \\
            }
        }
    }
     \\ \\

\end{supertabular}
}

\end{document}
